\begin{frame}[fragile,shrink=5]
  \frametitle{Marco Teórico: Ecuaciones de Navier-Stokes}
  
  Conservación de momento y continuidad en fluidos viscosos:
  
  \begin{align*}
    \rho \left( \frac{\partial \mathbf{u}}{\partial t} + \mathbf{u} \cdot \nabla \mathbf{u} \right) &= -\nabla p + \mu \nabla^2 \mathbf{u} + \mathbf{f} \\
    \nabla \cdot \mathbf{u} &= 0
  \end{align*}
  
  Donde: $\mathbf{u}$ = velocidad, $p$ = presión, $\rho$ = densidad, $\mathbf{f}$ = fuerzas externas (Coriolis).
  
  \begin{block}{Reto clave}
    El término convectivo $(\mathbf{u} \cdot \nabla \mathbf{u})$ introduce no linealidad que dificulta soluciones analíticas.
  \end{block}
  
  \note[item]{Tiempo estimado: 3 minutos}
\end{frame}

\begin{frame}[fragile]
  \frametitle{PINNs: Redes Neuronales Informadas por Física}
  
  \begin{block}{Concepto clave}
    Las PINNs codifican ecuaciones diferenciales parciales directamente en su función de pérdida, garantizando que las predicciones respeten leyes físicas.
  \end{block}
  
  \begin{itemize}
    \item Combinan datos observados con restricciones físicas
    \item No requieren volúmenes masivos de datos
    \item Conocimiento físico actúa como regularizador
    \item 30\% más precisión que NWP en casos de intensificación rápida
  \end{itemize}
  
  \note[item]{Tiempo estimado: 2 minutos}
\end{frame}
